\newpage
\section{Déroulement du projet}
\subsection{Réalisation du projet par étapes}
On a découpé le projet en plusieurs étapes successives :
\begin{enumerate}
\item Lecture et analyse des spécifications, puis conception de l'architecture MVC et des diagrammes UML ;
\item Mise en place des classes de base : plateau, cases, hiérarchie des pièces ;
\item Implémentation des règles de déplacement de chaque pièce, en gérant les colonnes neutres ;
\item Développement de la logique de validation des coups, de la détection d'échec et de la règle face à face ;
\item Développement de l'interface graphique avec Java Swing ;
\item Implémentation du bot et de ses trois niveaux de difficulté ;
\item Mise en place du système de logs avec Log4j ;
\item Écriture des tests unitaires JUnit ;
\item Rédaction du rapport.
\end{enumerate}
\subsection{Répartition des tâches, communication et intégration}
\paragraph{KETFI Ramdane}
S'est occupé du moteur de jeu et des classes de données : le plateau, la hiérarchie des pièces, la validation des coups, la détection d'échec, l'échec et mat, et la règle du face à face.
\paragraph{SOW Mamadou Saliou}
S'est occupé de l'interface graphique : la page d'accueil, le plateau de jeu, l'affichage des pièces capturées et l'historique des coups.
\paragraph{DJERMOUNE Thanina}
S'est occupée du bot et de ses trois niveaux de difficulté, ainsi que de la rédaction du rapport.
\paragraph{Communication et intégration}
Pour communiquer et partager notre travail, on a utilisé un groupe WhatsApp où on s'envoyait régulièrement les fichiers du projet sous forme de zip. On a aussi utilisé GitHub pour centraliser le code et éviter les conflits entre nos versions. Cela nous permettait de toujours avoir la dernière version du code et d'avancer ensemble malgré nos emplois du temps différents.
\paragraph{Difficultés rencontrées}
Au cours du projet, nous avons fait face à deux types de difficultés :
\begin{itemize}
\item \textbf{Difficultés organisationnelles :} Des divergences d'opinions ont parfois émergé au sein de l'équipe concernant certains choix techniques et de conception. Cependant, nous avons toujours privilégié le dialogue et la recherche d'un consensus, ce qui nous a permis de maintenir une cohésion solide tout au long du projet.
\item \textbf{Difficultés techniques :} La modélisation initiale du plateau et la hiérarchie des classes, notamment la gestion des colonnes neutres et de la rivière, ont nécessité une réflexion approfondie sur l'architecture du projet. L'intégration du bot au moteur de jeu a également été complexe, en particulier la conception d'un algorithme d'évaluation cohérent pour les niveaux Moyen et Difficile. Ces défis nous ont poussés à revoir certaines de nos décisions de conception, et nous en sommes sortis avec une meilleure compréhension du développement logiciel en équipe.
\end{itemize}